

\chapter{Supported boards}

\section{Introduction}

HobbySpark currently supports 30 different boards. Each board is identified
within HobbySpark by its corresponding class name. These class names are used
with the \code{set\_board} function when selecting the target board.

This chapter provides a quick hardware reference for every supported board.
The information includes the microcontroller, digital and analog I/O,
PWM capability, and commonly available communication interfaces.

The capabilities listed here describe the underlying hardware of the boards.
They should not be interpreted as a guarantee that every listed feature is
currently exposed through HobbySpark.

\section{AdafruitFeatherESP32S2}

\textbf{Description:} An Adafruit Feather board based on the ESP32-S2
microcontroller.

\begin{table}[H]
\centering
\caption{Adafruit Feather ESP32-S2 hardware reference}
\label{tab:board-feather-esp32s2}
\begin{tabular}{ll}
\toprule
\textbf{Property} & \textbf{Specification} \\
\midrule
Microcontroller & ESP32-S2 \\
Digital I/O & 10 exposed logic pins \\
Analog input & 6 pins \\
PWM & Yes \\
I\textsuperscript{2}C & Yes \\
SPI & Yes \\
UART & Yes \\
DAC & 2 pins \\
\bottomrule
\end{tabular}
\end{table}

The ESP32-S2 provides flexible peripheral multiplexing, allowing interfaces
such as PWM, UART, I\textsuperscript{2}C, and SPI to be assigned to suitable
GPIO pins. The board provides six analog-capable pins, with two also capable
of DAC output.

\newpage
\section{AdafruitFeatherM0}

\textbf{Description:} An Adafruit Feather board based on the SAMD21
microcontroller.

\begin{table}[H]
\centering
\caption{Adafruit Feather M0 hardware reference}
\label{tab:board-feather-m0}
\begin{tabular}{ll}
\toprule
\textbf{Property} & \textbf{Specification} \\
\midrule
Microcontroller & ATSAMD21G18 \\
Digital I/O & 20+ exposed I/O pins \\
Analog input & Up to 8 analog inputs \\
PWM & Yes \\
I\textsuperscript{2}C & Yes \\
SPI & Yes \\
UART & Yes \\
DAC & Yes \\
\bottomrule
\end{tabular}
\end{table}

The Feather M0 operates using 3.3 V logic. Its SAMD21 microcontroller
provides flexible timer and peripheral multiplexing. Nearly all of the
exposed logic pins can be used for PWM, although particular peripheral
assignments can impose pin-specific restrictions.

\section{ArduinoDue}

\textbf{Description:} An Arduino board based on the 32-bit SAM3X8E
microcontroller.

\begin{table}[H]
\centering
\caption{Arduino Due hardware reference}
\label{tab:board-arduino-due}
\begin{tabular}{ll}
\toprule
\textbf{Property} & \textbf{Specification} \\
\midrule
Microcontroller & SAM3X8E \\
Digital I/O & 54 \\
Analog input & 12 \\
PWM & 12 pins \\
I\textsuperscript{2}C & 2 interfaces \\
SPI & Yes \\
UART & 4 hardware UARTs \\
DAC & 2 \\
CAN & 2 interfaces \\
\bottomrule
\end{tabular}
\end{table}

The Arduino Due is based on a 32-bit ARM Cortex-M3 processor operating at
84 MHz. It provides substantially more I/O and communication capability
than many of the smaller AVR-based Arduino boards.

\newpage
\section{ArduinoGigaR1}

\textbf{Description:} An Arduino board designed for more advanced projects
and based on the STM32H747XI microcontroller.

\begin{table}[H]
\centering
\caption{Arduino GIGA R1 hardware reference}
\label{tab:board-giga-r1}
\begin{tabular}{ll}
\toprule
\textbf{Property} & \textbf{Specification} \\
\midrule
Microcontroller & STM32H747XI \\
Digital I/O & 76 \\
Analog input & 12 \\
PWM & Yes \\
I\textsuperscript{2}C & Yes \\
SPI & Yes \\
UART & Multiple hardware UARTs \\
DAC & Yes \\
CAN & Yes \\
\bottomrule
\end{tabular}
\end{table}

The GIGA R1 uses a high-performance STM32H747XI microcontroller and provides
a large number of general-purpose I/O pins together with multiple hardware
communication peripherals.

\section{ArduinoLeonardo}

\textbf{Description:} An Arduino board based on the ATmega32U4
microcontroller.

\begin{table}[H]
\centering
\caption{Arduino Leonardo hardware reference}
\label{tab:board-leonardo}
\begin{tabular}{ll}
\toprule
\textbf{Property} & \textbf{Specification} \\
\midrule
Microcontroller & ATmega32U4 \\
Digital I/O & 20 \\
Analog input & 12 \\
PWM & 7 pins \\
I\textsuperscript{2}C & Yes \\
SPI & Yes \\
UART & 1 hardware UART \\
USB & Native USB \\
\bottomrule
\end{tabular}
\end{table}

The ATmega32U4 provides native USB functionality, allowing the Leonardo to
appear as a USB device without requiring a separate USB-to-serial
microcontroller.

\newpage
\section{ArduinoLilyPad}

\textbf{Description:} A compact Arduino board designed for wearable and
e-textile projects.

\begin{table}[H]
\centering
\caption{Arduino LilyPad hardware reference}
\label{tab:board-lilypad}
\begin{tabular}{ll}
\toprule
\textbf{Property} & \textbf{Specification} \\
\midrule
Microcontroller & ATmega328P \\
Digital I/O & 14 \\
Analog input & 6 \\
PWM & 6 pins \\
I\textsuperscript{2}C & Yes \\
SPI & Yes \\
UART & Yes \\
\bottomrule
\end{tabular}
\end{table}

The LilyPad is designed for integration into wearable and textile-based
projects. Its electrical capabilities are broadly comparable to those of
ATmega328P-based Arduino boards.

\section{ArduinoMega2560}

\textbf{Description:} An Arduino board with a large number of digital and
analog pins, based on the ATmega2560 microcontroller.

\begin{table}[H]
\centering
\caption{Arduino Mega 2560 hardware reference}
\label{tab:board-mega2560}
\begin{tabular}{ll}
\toprule
\textbf{Property} & \textbf{Specification} \\
\midrule
Microcontroller & ATmega2560 \\
Digital I/O & 54 \\
Analog input & 16 \\
PWM & 15 pins \\
I\textsuperscript{2}C & Yes \\
SPI & Yes \\
UART & 4 hardware UARTs \\
\bottomrule
\end{tabular}
\end{table}

The Mega 2560 is particularly useful when a project requires a large number
of simultaneously connected digital devices.

\newpage
\section{ArduinoMicro}

\textbf{Description:} A compact Arduino board based on the ATmega32U4
microcontroller.

\begin{table}[H]
\centering
\caption{Arduino Micro hardware reference}
\label{tab:board-micro}
\begin{tabular}{ll}
\toprule
\textbf{Property} & \textbf{Specification} \\
\midrule
Microcontroller & ATmega32U4 \\
Digital I/O & 20 \\
Analog input & 12 \\
PWM & 7 pins \\
I\textsuperscript{2}C & Yes \\
SPI & Yes \\
UART & Yes \\
USB & Native USB \\
\bottomrule
\end{tabular}
\end{table}

\section{ArduinoMini}

\textbf{Description:} A small Arduino board designed for projects where a
compact form factor is useful.

\begin{table}[H]
\centering
\caption{Arduino Mini hardware reference}
\label{tab:board-mini}
\begin{tabular}{ll}
\toprule
\textbf{Property} & \textbf{Specification} \\
\midrule
Microcontroller & ATmega328P \\
Digital I/O & 14 \\
Analog input & 8 \\
PWM & 6 pins \\
I\textsuperscript{2}C & Yes \\
SPI & Yes \\
UART & Yes \\
\bottomrule
\end{tabular}
\end{table}

\section{ArduinoNano}

\textbf{Description:} A compact Arduino board based on the ATmega328P
microcontroller.

\begin{table}[H]
\centering
\caption{Arduino Nano hardware reference}
\label{tab:board-nano}
\begin{tabular}{ll}
\toprule
\textbf{Property} & \textbf{Specification} \\
\midrule
Microcontroller & ATmega328P \\
Digital I/O & 22 \\
Analog input & 8 \\
PWM & 6 pins \\
I\textsuperscript{2}C & Yes \\
SPI & Yes \\
UART & Yes \\
\bottomrule
\end{tabular}
\end{table}

\newpage
\section{ArduinoNanoESP32}

\textbf{Description:} An Arduino Nano board based on the ESP32-S3
microcontroller.

\begin{table}[H]
\centering
\caption{Arduino Nano ESP32 hardware reference}
\label{tab:board-nano-esp32}
\begin{tabular}{ll}
\toprule
\textbf{Property} & \textbf{Specification} \\
\midrule
Microcontroller & ESP32-S3 \\
Digital I/O & 30 GPIOs on the module \\
Analog input & Multiple ADC-capable GPIOs \\
PWM & Yes \\
I\textsuperscript{2}C & Yes \\
SPI & Yes \\
UART & Multiple UART interfaces \\
Wi-Fi & Yes \\
Bluetooth & Bluetooth LE \\
\bottomrule
\end{tabular}
\end{table}

The Nano ESP32 combines the compact Nano form factor with the wireless and
peripheral capabilities of the ESP32-S3.

\section{ArduinoNanoEvery}

\textbf{Description:} A compact Arduino Nano board based on the ATmega4809
microcontroller.

\begin{table}[H]
\centering
\caption{Arduino Nano Every hardware reference}
\label{tab:board-nano-every}
\begin{tabular}{ll}
\toprule
\textbf{Property} & \textbf{Specification} \\
\midrule
Microcontroller & ATmega4809 \\
Digital I/O & 20 \\
Analog input & 8 \\
PWM & Yes \\
I\textsuperscript{2}C & Yes \\
SPI & Yes \\
UART & Yes \\
\bottomrule
\end{tabular}
\end{table}

\section{ArduinoProMini}

\textbf{Description:} A small Arduino board designed for embedded projects
and based on the ATmega328P microcontroller.

\begin{table}[H]
\centering
\caption{Arduino Pro Mini hardware reference}
\label{tab:board-pro-mini}
\begin{tabular}{ll}
\toprule
\textbf{Property} & \textbf{Specification} \\
\midrule
Microcontroller & ATmega328P \\
Digital I/O & 14 \\
Analog input & 8 \\
PWM & 6 pins \\
I\textsuperscript{2}C & Yes \\
SPI & Yes \\
UART & Yes \\
\bottomrule
\end{tabular}
\end{table}

\section{ArduinoUnoR3}

\textbf{Description:} The third revision of the Arduino Uno, based on the
ATmega328P microcontroller.

\begin{table}[H]
\centering
\caption{Arduino Uno R3 hardware reference}
\label{tab:board-uno-r3}
\begin{tabular}{ll}
\toprule
\textbf{Property} & \textbf{Specification} \\
\midrule
Microcontroller & ATmega328P \\
Digital I/O & 14 \\
Analog input & 6 \\
PWM & 6 pins \\
I\textsuperscript{2}C & Yes \\
SPI & Yes \\
UART & 1 hardware UART \\
\bottomrule
\end{tabular}
\end{table}

The Uno R3 is one of the most widely used Arduino boards and is based on
the 8-bit ATmega328P microcontroller.

\section{ArduinoUnoWiFiRev2}

\textbf{Description:} An Arduino Uno board with Wi-Fi connectivity, based
on the ATmega4809 microcontroller.

\begin{table}[H]
\centering
\caption{Arduino Uno WiFi Rev2 hardware reference}
\label{tab:board-uno-wifi}
\begin{tabular}{ll}
\toprule
\textbf{Property} & \textbf{Specification} \\
\midrule
Microcontroller & ATmega4809 \\
Digital I/O & 20 \\
Analog input & 6 \\
PWM & 5 pins \\
I\textsuperscript{2}C & Yes \\
SPI & Yes \\
UART & Yes \\
Wi-Fi & Yes \\
Bluetooth LE & Yes \\
\bottomrule
\end{tabular}
\end{table}

\newpage
\section{ArduinoZero}

\textbf{Description:} An Arduino board based on the SAMD21 microcontroller.

\begin{table}[H]
\centering
\caption{Arduino Zero hardware reference}
\label{tab:board-zero}
\begin{tabular}{ll}
\toprule
\textbf{Property} & \textbf{Specification} \\
\midrule
Microcontroller & ATSAMD21G18 \\
Digital I/O & 20 \\
Analog input & 6 \\
PWM & Yes \\
I\textsuperscript{2}C & Yes \\
SPI & Yes \\
UART & Yes \\
DAC & 1 \\
\bottomrule
\end{tabular}
\end{table}

\section{ESP32C3}

\textbf{Description:} An ESP32 development board based on the ESP32-C3
microcontroller.

\begin{table}[H]
\centering
\caption{ESP32-C3 hardware reference}
\label{tab:board-esp32-c3}
\begin{tabular}{ll}
\toprule
\textbf{Property} & \textbf{Specification} \\
\midrule
Microcontroller & ESP32-C3 \\
Digital I/O & Depends on board variant \\
Analog input & ADC-capable GPIOs \\
PWM & Yes \\
I\textsuperscript{2}C & Yes \\
SPI & Yes \\
UART & Multiple UART capability \\
Wi-Fi & Yes \\
Bluetooth & Bluetooth LE \\
\bottomrule
\end{tabular}
\end{table}

\section{ESP32DevKit}

\textbf{Description:} An ESP32 development board intended for
general-purpose ESP32 projects.

\begin{table}[H]
\centering
\caption{ESP32 DevKit hardware reference}
\label{tab:board-esp32-devkit}
\begin{tabular}{ll}
\toprule
\textbf{Property} & \textbf{Specification} \\
\midrule
Microcontroller & ESP32 \\
Digital I/O & Up to 34 GPIOs on the ESP32 \\
Analog input & ADC-capable GPIOs \\
PWM & Yes \\
I\textsuperscript{2}C & Yes \\
SPI & Yes \\
UART & 3 UART interfaces \\
DAC & 2 channels \\
Wi-Fi & Yes \\
Bluetooth & Bluetooth Classic + BLE \\
\bottomrule
\end{tabular}
\end{table}

The exact number of exposed pins depends on the particular ESP32 DevKit
variant. The ESP32 supports peripheral multiplexing, allowing many
interfaces to be assigned to different GPIOs.

\section{ESP32S2}

\textbf{Description:} An ESP32 development board based on the ESP32-S2
microcontroller.

\begin{table}[H]
\centering
\caption{ESP32-S2 hardware reference}
\label{tab:board-esp32-s2}
\begin{tabular}{ll}
\toprule
\textbf{Property} & \textbf{Specification} \\
\midrule
Microcontroller & ESP32-S2 \\
Digital I/O & Depends on board variant \\
Analog input & ADC-capable GPIOs \\
PWM & Yes \\
I\textsuperscript{2}C & Yes \\
SPI & Yes \\
UART & Yes \\
DAC & 2 channels \\
Wi-Fi & Yes \\
Bluetooth & No Bluetooth radio \\
\bottomrule
\end{tabular}
\end{table}

\section{ESP32S3}

\textbf{Description:} An ESP32 development board based on the ESP32-S3
microcontroller.

\begin{table}[H]
\centering
\caption{ESP32-S3 hardware reference}
\label{tab:board-esp32-s3}
\begin{tabular}{ll}
\toprule
\textbf{Property} & \textbf{Specification} \\
\midrule
Microcontroller & ESP32-S3 \\
Digital I/O & Depends on board variant \\
Analog input & ADC-capable GPIOs \\
PWM & Yes \\
I\textsuperscript{2}C & Yes \\
SPI & Yes \\
UART & 3 UART interfaces \\
Wi-Fi & Yes \\
Bluetooth & Bluetooth LE \\
\bottomrule
\end{tabular}
\end{table}

\newpage
\section{ESP8266D1Mini}

\textbf{Description:} A compact development board based on the ESP8266
microcontroller.

\begin{table}[H]
\centering
\caption{ESP8266 D1 Mini hardware reference}
\label{tab:board-esp8266-d1-mini}
\begin{tabular}{ll}
\toprule
\textbf{Property} & \textbf{Specification} \\
\midrule
Microcontroller & ESP8266EX \\
Digital I/O & 11 GPIOs on the ESP8266 \\
Analog input & 1 ADC input \\
PWM & Yes \\
I\textsuperscript{2}C & Software/peripheral support \\
SPI & Yes \\
UART & Yes \\
Wi-Fi & Yes \\
\bottomrule
\end{tabular}
\end{table}

\section{ESP8266NodeMCU}

\textbf{Description:} A NodeMCU development board based on the ESP8266
microcontroller.

\begin{table}[H]
\centering
\caption{ESP8266 NodeMCU hardware reference}
\label{tab:board-esp8266-nodemcu}
\begin{tabular}{ll}
\toprule
\textbf{Property} & \textbf{Specification} \\
\midrule
Microcontroller & ESP8266EX \\
Digital I/O & Up to 17 GPIOs on the MCU \\
Analog input & 1 ADC input \\
PWM & Yes \\
I\textsuperscript{2}C & Yes \\
SPI & Yes \\
UART & Yes \\
Wi-Fi & Yes \\
\bottomrule
\end{tabular}
\end{table}

\section{RaspberryPiPico}

\textbf{Description:} A compact Raspberry Pi microcontroller board based
on the RP2040 microcontroller.

\begin{table}[H]
\centering
\caption{Raspberry Pi Pico hardware reference}
\label{tab:board-pico}
\begin{tabular}{ll}
\toprule
\textbf{Property} & \textbf{Specification} \\
\midrule
Microcontroller & RP2040 \\
Digital I/O & 26 exposed GPIOs \\
Analog input & 3 \\
PWM & 16 channels \\
I\textsuperscript{2}C & 2 controllers \\
SPI & 2 controllers \\
UART & 2 controllers \\
\bottomrule
\end{tabular}
\end{table}

The RP2040 provides flexible programmable I/O and multiple hardware
communication interfaces.

\newpage
\section{RaspberryPiPicoW}

\textbf{Description:} A Raspberry Pi Pico board with wireless connectivity
based on the RP2040 microcontroller.

\begin{table}[H]
\centering
\caption{Raspberry Pi Pico W hardware reference}
\label{tab:board-pico-w}
\begin{tabular}{ll}
\toprule
\textbf{Property} & \textbf{Specification} \\
\midrule
Microcontroller & RP2040 \\
Digital I/O & 26 exposed GPIOs \\
Analog input & 3 \\
PWM & 16 channels \\
I\textsuperscript{2}C & 2 controllers \\
SPI & 2 controllers \\
UART & 2 controllers \\
Wi-Fi & Yes \\
Bluetooth & Bluetooth LE \\
\bottomrule
\end{tabular}
\end{table}

The Pico W provides essentially the same RP2040-based I/O architecture as
the Pico while adding wireless connectivity.

\section{STM32BlackPill}

\textbf{Description:} A compact development board based on an STM32
microcontroller.

\begin{table}[H]
\centering
\caption{STM32 Black Pill hardware reference}
\label{tab:board-black-pill}
\begin{tabular}{ll}
\toprule
\textbf{Property} & \textbf{Specification} \\
\midrule
Microcontroller & STM32 family; varies by Black Pill variant \\
Digital I/O & Depends on variant \\
Analog input & Depends on variant \\
PWM & Yes \\
I\textsuperscript{2}C & Yes \\
SPI & Yes \\
UART & Multiple hardware UARTs \\
\bottomrule
\end{tabular}
\end{table}

The name ``Black Pill'' is used for several STM32 development boards.
Consequently, the exact processor and number of exposed pins depend on the
specific hardware variant.

\newpage
\section{STM32BluePill}

\textbf{Description:} A compact development board based on an STM32
microcontroller.

\begin{table}[H]
\centering
\caption{STM32 Blue Pill hardware reference}
\label{tab:board-blue-pill}
\begin{tabular}{ll}
\toprule
\textbf{Property} & \textbf{Specification} \\
\midrule
Microcontroller & STM32F103C8T6 (common variant) \\
Digital I/O & Up to 37 GPIOs \\
Analog input & Up to 10 ADC channels \\
PWM & Yes \\
I\textsuperscript{2}C & 2 interfaces \\
SPI & 2 interfaces \\
UART & 3 USART interfaces \\
\bottomrule
\end{tabular}
\end{table}

The Blue Pill is commonly based on the STM32F103C8T6 microcontroller and
provides a relatively large number of GPIO pins in a compact form factor.

\section{SeeeduinoXIAO}

\textbf{Description:} A compact Seeeduino board designed for small
embedded projects.

\begin{table}[H]
\centering
\caption{Seeeduino XIAO hardware reference}
\label{tab:board-xiao}
\begin{tabular}{ll}
\toprule
\textbf{Property} & \textbf{Specification} \\
\midrule
Microcontroller & ATSAMD21G18 \\
Digital I/O & 11 \\
Analog input & 11-capable GPIOs \\
PWM & Yes \\
I\textsuperscript{2}C & Yes \\
SPI & Yes \\
UART & Yes \\
DAC & Yes \\
\bottomrule
\end{tabular}
\end{table}

The XIAO provides a very small form factor while retaining the peripheral
capabilities of the SAMD21 microcontroller.

\newpage
\section{SeeeduinoXIAOESP32C3}

\textbf{Description:} A Seeeduino XIAO board based on the ESP32-C3
microcontroller.

\begin{table}[H]
\centering
\caption{Seeed XIAO ESP32-C3 hardware reference}
\label{tab:board-xiao-esp32-c3}
\begin{tabular}{ll}
\toprule
\textbf{Property} & \textbf{Specification} \\
\midrule
Microcontroller & ESP32-C3 \\
Digital I/O & 11 \\
Analog input & 4 \\
PWM & Yes \\
I\textsuperscript{2}C & Yes \\
SPI & Yes \\
UART & Yes \\
Wi-Fi & Yes \\
Bluetooth & Bluetooth LE \\
\bottomrule
\end{tabular}
\end{table}

The XIAO ESP32-C3 combines the small XIAO form factor with the wireless
capabilities of the ESP32-C3.
\section{Teensy40}

\textbf{Description:} A compact Teensy development board based on the
NXP i.MX RT1062 microcontroller.

\begin{table}[H]
\centering
\caption{Teensy 4.0 hardware reference}
\label{tab:board-teensy40}
\begin{tabular}{ll}
\toprule
\textbf{Property} & \textbf{Specification} \\
\midrule
Microcontroller & NXP i.MX RT1062 \\
Digital I/O & 40 \\
Analog input & 14 \\
PWM & Yes \\
I\textsuperscript{2}C & Multiple interfaces \\
SPI & Multiple interfaces \\
UART & Multiple hardware serial ports \\
CAN & Yes \\
\bottomrule
\end{tabular}
\end{table}

The Teensy 4.0 uses a high-performance ARM Cortex-M7 processor and is
designed for applications requiring substantially greater processing
performance than typical 8-bit Arduino boards.

\newpage
\section{Teensy41}

\textbf{Description:} A Teensy development board based on the NXP i.MX RT1062
microcontroller with additional memory and connectivity compared with the
Teensy 4.0.

\begin{table}[H]
\centering
\caption{Teensy 4.1 hardware reference}
\label{tab:board-teensy41}
\begin{tabular}{ll}
\toprule
\textbf{Property} & \textbf{Specification} \\
\midrule
Microcontroller & NXP i.MX RT1062 \\
Digital I/O & 55 \\
Analog input & 18 \\
PWM & Yes \\
I\textsuperscript{2}C & Multiple interfaces \\
SPI & Multiple interfaces \\
UART & Multiple hardware serial ports \\
CAN & Yes \\
Ethernet & Yes \\
\bottomrule
\end{tabular}
\end{table}

The Teensy 4.1 provides the same high-performance processor as the Teensy
4.0 while offering substantially more memory, I/O, and expansion capability.

\newpage
\section{Summary}

Table~\ref{tab:all-supported-boards} provides a compact summary of all
boards currently supported by HobbySpark.

\begin{table}[H]
\centering
\caption{HobbySpark supported boards}
\label{tab:all-supported-boards}
\begin{tabular}{ll}
\toprule
\textbf{HobbySpark Board Class} & \textbf{Microcontroller} \\
\midrule
AdafruitFeatherESP32S2 & ESP32-S2 \\
AdafruitFeatherM0 & ATSAMD21 \\
ArduinoDue & SAM3X8E \\
ArduinoGigaR1 & STM32H747XI \\
ArduinoLeonardo & ATmega32U4 \\
ArduinoLilyPad & ATmega328P \\
ArduinoMega2560 & ATmega2560 \\
ArduinoMicro & ATmega32U4 \\
ArduinoMini & ATmega328P \\
ArduinoNano & ATmega328P \\
ArduinoNanoESP32 & ESP32-S3 \\
ArduinoNanoEvery & ATmega4809 \\
ArduinoProMini & ATmega328P \\
ArduinoUnoR3 & ATmega328P \\
ArduinoUnoWiFiRev2 & ATmega4809 \\
ArduinoZero & ATSAMD21 \\
ESP32C3 & ESP32-C3 \\
ESP32DevKit & ESP32 \\
ESP32S2 & ESP32-S2 \\
ESP32S3 & ESP32-S3 \\
ESP8266D1Mini & ESP8266 \\
ESP8266NodeMCU & ESP8266 \\
RaspberryPiPico & RP2040 \\
RaspberryPiPicoW & RP2040 \\
STM32BlackPill & STM32 family \\
STM32BluePill & STM32F103C8T6 \\
SeeeduinoXIAO & ATSAMD21 \\
SeeeduinoXIAOESP32C3 & ESP32-C3 \\
Teensy40 & i.MX RT1062 \\
Teensy41 & i.MX RT1062 \\
\bottomrule
\end{tabular}
\end{table}